\clearpage\begingroup
\newfontfamily\ReichenbachSymbolFont{FreeSerif.otf}
\SourceLeaf{322}\NoSourceHead
\begin{center}
{\large [BEOBACHTUNGEN AM\\REICHENBACHSCHEN MERIDIANKREISE.]}\par\medskip
BRIEFWECHSEL.\par\medskip
\rule{\textwidth}{.3pt}\par
{\small [\textsc{Wilhelm Olbers}, Sein Leben und seine Werke, II1, Berlin 1900 und II2, Berlin 1909. Briefwechsel zwischen \textsc{Gauss} und \textsc{Bessel}, Leipzig 1880. Briefwechsel zwischen \textsc{Olbers} und \textsc{Bessel}, 2 Bände, Leipzig 1852.]}\par
\rule{\textwidth}{.3pt}\par\medskip
\textsc{Gauss} an \textsc{Olbers}[, Göttingen, Mai 1819].
\end{center}

\noindent\hspace*{\parindent}. . . . . . Ueber die Constantenfehler der Höhenmessungs-Instrumente sind noch von Niemand Untersuchungen angestellt, auch von \textsc{Bessel} nicht. Meine Meinung ist, dass jedes Instrument hiebei ein Individuum ist, und dass wir das Wahre noch gar nicht kennen. Jedes Instrument wird seinem Bau nach, sowie wegen der \textls[120]{Biegungen} der Theile Fehler der gemessenen Zenithdistanzen von der Form \(\alpha\sin z+\beta\cos z\) haben; der zweite Theil fällt weg bei Instrumenten, die umgewandt werden können, also nicht bei \textsc{Pond}s Kreise, aller Wahrscheinlichkeit nach ist er aber unmerklich. Allein der erste Theil kann meiner Meinung nach durch alles Aequilibriren nicht mit Gewissheit weggeschafft werden, und ich sehe durchaus kein Mittel, ihn zu erforschen, als durch den {\ReichenbachSymbolFont ☿}-Horizont[\begingroup\renewcommand{\thefootnote}{*)}\footnotemark\renewcommand{\thefootnote}{[*)}\footnotetext{{\ReichenbachSymbolFont ☿} \(=\) Quecksilber.]}\endgroup]. Leider sind diese Beob. sehr mühsam und ich habe bisher nur erst ein Paar gemacht, die zeigen, dass \(\alpha\) bei \textsc{Repsold}s Kreise nur sehr klein sein kann. Wäre \textsc{Bessel}s [\textsc{Cary}scher] Kreis fehlerfrei, so müsste, nach den Resultaten zu schliessen, \(\alpha\) bei \textsc{Pond}s Kreise
\clearpage\SourceLeaf{323}\HistoricalHead{318}{L}{PRAKTISCHE UND SPHÄRISCHE ASTRONOMIE. BRIEFWECHSEL.}
\noindent etwa \(=-3\overset{\prime\prime}{,}3\) u[nd] bei \textsc{Repsold}s Kreise \(=-4\overset{\prime\prime}{,}8\) sein, was anzunehmen mir schwer wird; eher möchte ich zugeben, dass \textsc{Pond}s Kreis fehlerfrei sei und \(\alpha=-1\overset{\prime\prime}{,}5\) bei mir u[nd] \(=+3\overset{\prime\prime}{,}3\) bei \textsc{Bessel}. . . . . . .

\begin{center}\textsc{Bessel} an \textsc{Gauss}, 20. März 1820.\end{center}
\noindent\hspace*{\parindent}\{. . . . . . Was das neue Instr[ument, den Königsberger \textsc{Reichenbach}schen Meridiankreis] anlangt, so wird seine Güte hauptsächlich von der Unveränderlichkeit seiner Collimationslinie abhängen, worüber mir, des Wetters und einiger Veränderungen des Fadennetzes wegen, bis jetzt die Erfahrungen fehlen. Zu Anfang werde ich, um über diesen Punkt aufs Reine zu kommen, das Instr[ument] sehr häufig umlegen; später und wenn erst die Biegung des Fernrohrs untersucht sein wird, wird es zweckmässiger sein, die Kreise lange unberührt zu lassen. — Inzwischen zweifle ich nicht daran, dass fortgesetzte Beob[achtungen] manche Winke an die Hand geben werden, wie das Instrument recht zweckmässig angewandt werden kann. — Sie haben gewiss Recht, wenn Sie den Gebrauch der starken Vergrösserungen anrathen; allein die Uebung und Gewohnheit macht auch hier sehr viel, und wenn es sich zeigen sollte, dass mit dem zweiten Oculare (dem nächst schwächsten) die Beob[achtung] eben so genau gemacht werden kann, so würde dieses den Vorzug haben, dass man dann das Fernrohr gar nicht berühren dürfte, indem dieses Ocular alle Fäden hinlänglich übersehen lässt. — Wenn ich früher mit 44m[aliger] Vergr[össerung] und 3 starken Metallfäden einen ziemlich kleinen Beobachtungsfehler erhielt (\(0\overset{\prime\prime}{,}1022\)), so ist zu erwarten, dass die Vergrösserung, wovon ich rede, 5 feine Fäden und ein lichtstärkeres Fernrohr ihn auf die Hälfte reduciren werden, in welchem Falle ich die Annehmlichkeit, die das erwähnte Ocular darbietet, doch dem von der stärkeren Vergr[össerung] vielleicht zu erwartenden Vortheile \textls[120]{gleich} achten würde. Sobald anhaltend heiteres Wetter einfällt, will ich mit den verschiedenen Ocularen Versuche machen. . . . . . .

. . . . . . Vier Beob[achtungen] des Polarsterns, wovon die 2 ersten und die 2 letzten in entgegengesetzter Lage gemacht sind, geben mir die Polhöhe \(54^\circ42'50''\,79\). Später habe ich einen fehlgeschlagenen Versuch gemacht, dieses Sterns Z[enith]d[istanz] aus \textls[120]{Oel} zu nehmen; das Gefäss war etwa 8 Zoll breit und 16 Zoll lang, allein doch war die Oberfl[äche] nicht eben,
\clearpage\SourceLeaf{324}\HistoricalHead{319}{R}{BEOBACHTUNGEN AM REICHENBACHSCHEN MERIDIANKREISE.}
\noindent wie das undeutliche Bild des Sterns zeigte. Ich werde nun Wasser versuchen, was wohl besser gelingen wird. . . . . . .\} [\begingroup\renewcommand{\thefootnote}{*)}\footnotemark\renewcommand{\thefootnote}{[*)}\footnotetext{Vermuthlich war die Krümmung des Quecksilberhorizonts infolge der Kapillarität der Grund, warum weder \textsc{Gauss} noch \textsc{Bessel} diesen sogleich regelmässig benutzten. Vgl. die auf Seite 342 abgedruckte Briefstelle.]}\endgroup]

\begin{center}\textsc{Olbers} an \textsc{Bessel}, 20. April 1820.\end{center}
\noindent\hspace*{\parindent}\{. . . . . . Weder von \textsc{Pond}s noch von \textsc{Gauss}’ Erfolge der Versuche, auf einem Quecksilberhorizont zu beobachten, ist mir bisher etwas Näheres bekannt geworden. . . . . . .\}

\begin{center}\textsc{Bessel} an \textsc{Gauss}, 30. April 1820.\end{center}
\noindent\hspace*{\parindent}\{. . . . . . Ueber die Polhöhe habe ich [mit dem neuen \textsc{Reichenbach}schen Meridiankreise] nur wenige Resultate, weil das Wetter nur selten ruhig genug ist; diese wenigen stimmen aber gut überein und ich theile sie Ihnen mit, da Sie selbst am besten wissen, welcher Werth solchen einseitigen Bestimmungen gebührt. Es sind Beobachtungen der doppelten Höhe des Polarsterns von einem \textls[120]{Wasser}horizonte, die mit der Declination aus meinen Tafeln berechnet wurden. . . . . . .

. . . . . . Hiernach wird die Polhöhe nahe an \([54^\circ42']50''\) kommen, wenngleich noch viel fehlt, ehe man dieses bis auf \(0\overset{\prime\prime}{,}1\) oder \(0\overset{\prime\prime}{,}2\) behaupten kann. Directe Beobachtungen würden etwa \(51\overset{\prime\prime}{,}6\) geben oder die Biegung im Horizonte etwa \(=3''\). — Hätte man die strenge Prüfung durch Beob[achtungen] von einem Horizonte nach meinem früheren Vorschlage bei den Wiederholungskreisen angewandt, so würde man vielleicht auch hier stets dieselben absoluten Höhen gefunden haben; denn diese Instrumente scheinen mir noch \textls[120]{weit} gefährlicher zu sein als unsere einfacheren Meridiankreise, zumal da Flexion der Kreise und der von \textsc{Bohnenberger} aufgefundene Fehler[\begingroup\renewcommand{\thefootnote}{**)}\footnotemark\renewcommand{\thefootnote}{[**)}\footnotetext{Wirkung der Schwere auf die Achse: \textsc{Bohnenberger}, \textit{Zusatz zu dem} — (Werke VI, S. 393 abgedruckten) — \textit{Schreiben des Herrn Hofrath Gauss}, Zeitschrift für Astronomie, Bd. 4, S. 141, Juli-August 1817.]}\endgroup] hinzukommen. Mein \textsc{Cary}scher Kreis ist aber so gebaut, dass an eine Flexion gar nicht zu denken ist; — wenn der neue [\textsc{Reichenbach}sche] Kr[eis] dieselbe Polhöhe gibt und sonach der Grund der Zweifel entweder geschwächt wird oder gar verschwindet, so wird es unnöthig sein, die lange projectirte Beobachtungsreihe über den Polarstern aus einem Horizonte wirklich zu machen. . . . . . . .\}
\clearpage\endgroup
